%! Function Operations — Unit 1, Lesson 1.2 — Warm-Up
% Exactly ONE page, blank and keyed (LESSON_SHAPE.md, one_page). Four prerequisite items:
% 1 evaluates a function at a negative input (1.1's habit, needed for combining at a point);
% 2 subtracts a whole expression, the move (f-g) lives or dies on; 3 multiplies two binomials,
% which (fg) needs; 4 solves a "bottom equals zero" equation — the exact skill the quotient
% domain runs on.
\documentclass[10pt]{article}
\usepackage{precalculus-article}
\usepackage{precalculus-boxes}
\newcommand{\TallMath}[1]{$\displaystyle #1\rule[-1.4em]{0pt}{3.2em}$}

\begin{document}

\pageheader{Unit 1, Lesson 1.2}{Warm-Up}

\vspace{0.12in}

\begin{enumerate}[itemsep=10pt]

\item Let $f(x) = 2x - 7$. Find $f(-3)$.

\begin{work}
  f(-3) &= 2(-3) - 7 \\
        &= -6 - 7 = -13
\end{work}

$f(-3) = $ \blank{3cm}

\item Simplify by removing the parentheses first: \quad $(4x + 9) - (x - 6)$

\begin{work}
  (4x + 9) - (x - 6) &= 4x + 9 - x + 6 \\
                     &= 3x + 15
\end{work}

Simplified: \blank{3cm}

\item Multiply and simplify: \quad $(x + 4)(x - 2)$

\begin{work}
  (x + 4)(x - 2) &= x^2 - 2x + 4x - 8 \\
                 &= x^2 + 2x - 8
\end{work}

Product: \blank{3cm}

\item Solve each equation. These ask, ``what input makes this expression equal to zero?''

\begin{enumerate}[label=(\alph*), itemsep=8pt]
  \item $x - 3 = 0$ \qquad $x = $ \blank{2.4cm}
  \item $2x - 10 = 0$ \qquad $x = $ \blank{2.4cm}
\end{enumerate}

\end{enumerate}

\end{document}
